% !TeX program = xelatex
\documentclass[10pt,a4paper]{article}

% ==================================================
% DOCUMENT METADATA
% ==================================================

% ==================================================
% DOCUMENT METADATA
% PROJECT ASSIGNMENT 2 (PA2)
% ==================================================
% Escape &, %, _, and # as \&, \%, \_, and \#.

% Project Information
\newcommand{\ProjectName}{MEDICATION SAFETY ASSISTANT}
\newcommand{\GroupName}{Medication Safety Assistant Team}
\newcommand{\GroupNumber}{Group 02}
\newcommand{\OrganizationName}{University of Science, VNU-HCM}

% Document Information
\newcommand{\DocumentTitle}{Use-Case Specification}
\newcommand{\DocumentVersion}{1.1}
\newcommand{\DocumentDate}{1 November 2026}
\newcommand{\DocumentIdentifier}{PA2-Group02-UseCaseSpecification}

% Document Authors
\newcommand{\DocumentAuthors}{%
    Cao Hoàng Long, Hoàng Yến Trang, Lỗ Chí Thành%
}

% Document Classification
\newcommand{\ConfidentialityText}{Course Project}


% ==================================================
% PACKAGES
% ==================================================
\usepackage{fontspec}
\setmainfont{Times New Roman}

\usepackage[
    a4paper,
    top=1.6cm,
    bottom=1.6cm,
    left=1.2cm,
    right=1.2cm,
    headheight=15pt
]{geometry}

\usepackage{setspace}
\setstretch{1.05}

\usepackage{xcolor}
\usepackage{array,longtable,tabularx}
\usepackage{graphicx}
\usepackage{float}
\usepackage{caption}
\usepackage{fancyhdr}
\usepackage{pdflscape}
\usepackage{titlesec}
\usepackage[unicode,hidelinks]{hyperref}

% build.ps1 converts the original SVG diagrams to vector PDF first.
\graphicspath{{sections/images/}}

% ==================================================
% COLORS AND CUSTOM COMMANDS
% ==================================================
\definecolor{instructionblue}{RGB}{0,92,170}

\newcommand{\instruction}[1]{%
    \par
    \textcolor{instructionblue}{\itshape[#1]}
    \par
}

\newcommand{\field}[1]{%
    \textcolor{instructionblue}{\itshape<#1>}
}

% Table column types
\newcolumntype{L}[1]{%
    >{\raggedright\arraybackslash}p{#1}
}

\newcolumntype{Y}{%
    >{\raggedright\arraybackslash}X
}

% ==================================================
% DOCUMENT FORMATTING
% ==================================================
\renewcommand{\arraystretch}{1.25}
\setlength{\tabcolsep}{3pt}
\emergencystretch=3em

\setlength{\parindent}{0pt}
\setlength{\parskip}{5pt}
\setlength{\LTpre}{0pt}
\setlength{\LTpost}{0pt}

\setcounter{secnumdepth}{2}
\setcounter{tocdepth}{2}

\titleformat{\section}
    {\bfseries\large}
    {\thesection.}
    {0.5em}
    {}

\titleformat{\subsection}
    {\bfseries}
    {\thesubsection}
    {0.5em}
    {}

% ==================================================
% HEADER AND FOOTER
% ==================================================
\pagestyle{fancy}
\fancyhf{}

\fancyhead[L]{\ProjectName}
\fancyhead[R]{Version \DocumentVersion}

\fancyfoot[L]{%
    \OrganizationName\ - \ConfidentialityText
}
\fancyfoot[R]{Page \thepage}

\renewcommand{\headrulewidth}{0.4pt}
\renewcommand{\footrulewidth}{0.4pt}

% ==================================================
% PDF METADATA
% ==================================================
\hypersetup{
    pdftitle={\DocumentTitle\ (\ProjectName)},
    pdfauthor={\DocumentAuthors},
    pdfsubject={PA2 Use-Case Specification}
}

% ==================================================
% SECTION LOADER
% ==================================================
\newcommand{\loadsection}[2]{%
    \IfFileExists{sections/#1.tex}{%
        \input{sections/#1}%
    }{%
        \section{#2}
        \textit{Content to be completed.}
    }%
}

% ==================================================
% DOCUMENT
% ==================================================
\begin{document}
\hypersetup{pageanchor=false}

% ==================================================
% TITLE PAGE
% ==================================================
\begin{titlepage}
    \centering

    \vspace*{2.2cm}

    {\Large\bfseries
        \ProjectName\par
    }

    \vspace{2cm}

    {\Huge\bfseries
        \DocumentTitle\par
    }

    \vspace{1.95cm}

    {\large
        Version \DocumentVersion\par
    }

    \vfill

    \begin{tabularx}{0.9\textwidth}{
        @{} L{3.5cm} Y @{}
    }
        \textbf{Document identifier:}
            & \DocumentIdentifier\\[6pt]

        \textbf{Prepared by:}
            & \DocumentAuthors\\[6pt]

        \textbf{Group:}
            & \GroupNumber\\[6pt]

        \textbf{Organization:}
            & \OrganizationName\\[6pt]

        \textbf{Date:}
            & \DocumentDate\\
    \end{tabularx}

    \vfill
\end{titlepage}
\hypersetup{pageanchor=true}

% ==================================================
% REVISION HISTORY
% ==================================================
\section*{Revision History}
\addcontentsline{toc}{section}{Revision History}

\begin{tabularx}{\textwidth}{
    |L{2.4cm}|L{1.7cm}|Y|L{3cm}|
}
\hline
\textbf{Date}
    & \textbf{Version}
    & \textbf{Description}
    & \textbf{Author}\\
\hline
1 November 2026
    & 1.0
    & Initial PA2 Use-Case Specification
    & Group 02\\
\hline
\textcolor{instructionblue}{1 November 2026}
    & \textcolor{instructionblue}{1.1}
    & \textcolor{instructionblue}{Updated portrait A4 layout and margins. UI prototype documentation is maintained as a separate report.}
    & \textcolor{instructionblue}{Group 02}\\
\hline
\end{tabularx}

% ==================================================
% TABLE OF CONTENTS
% ==================================================
\clearpage
\begingroup
\small
\setstretch{1.0}
\setlength{\parskip}{0pt}
\tableofcontents
\endgroup
\clearpage

% ==================================================
% MAIN CONTENT
% ==================================================

% Models: sections/01-use-case-model.tex
\loadsection{03-use-cases-combined}{Use-case Models and Specifications}

\clearpage

% Specifications: sections/02-use-case-specifications.tex
% The combined section includes the model followed by its specification.
% UI report is compiled separately from ../ui-prototype-report/main.tex.

% ==================================================
% END DOCUMENT
% ==================================================
\end{document}
